\section{Attention Timeline：Transformer 从何而来}

Transformer 的历史不能写成“2017 年突然发明 self-attention”。它继承了 sequence modeling、encoder-decoder、alignment 与 attention 的长期积累。时间线最重要的教学作用，是把 Transformer 看成对信息路径的重新设计：从递归压缩，转向按查询读取，再转向序列内部的直接交互。

\subsection{从 recurrent memory 到 attention}

\term{RNN}（recurrent neural network，循环神经网络）按时间递归更新隐藏状态：
\[
h_t = f(x_t,h_{t-1}),
\]
其中 $x_t$ 是第 $t$ 个输入，$h_t$ 是截至当前位置的状态，$f$ 通常由带参数的非线性单元实现。\term{LSTM} 与 \term{GRU} 是带门控的 RNN 变体，用显式 gate 缓解梯度消失并控制记忆保留。它们可以编码顺序，但信息必须沿长链逐步传播，训练也难以在时间维完全并行。

\lecturefigure{slide-05-attention-timeline.jpg}{Attention、self-attention 与预训练模型位于连续演化的时间线上。}{官方视频 03:10；Vaswani et al. 2017。}

\begin{knowledgebox}{读图：时间线的单位是信息路径}
最早的 sequence model 把历史压入 recurrent state；2014--2015 年 attention 允许 decoder 对多个 encoder state 加权读取；2017 年 Transformer 进一步让 token 之间直接 self-attend；之后的 BERT、GPT 等预训练模型把同一架构扩展为通用表示或生成接口。关键趋势不是模型名字越来越多，而是监督信号、计算并行性与可复用性同时扩大。
\end{knowledgebox}

\teachervoice{讲者把“Attention Is All You Need”称为简单 attention idea 的关键转折，但也先展示 prehistory。课堂提醒：一个成功架构往往把已有机制重新组合到更适合硬件、数据和训练目标的接口中，而不是从零创造全部组件。}

\subsection{Seq2seq 的 fixed-vector bottleneck}

\term{seq2seq}（sequence-to-sequence）模型把输入序列编码为表示，再由 decoder 逐步生成输出。早期做法常把整个源序列压成单个 \term{context vector}（上下文向量）$c=h_n$。当输入很长时，所有词义、顺序和依赖都竞争有限容量；decoder 还必须从同一 $c$ 中恢复每一步需要的信息。

\lecturefigure{slide-06-prehistory.jpg}{RNN、seq2seq、LSTM 与 GRU 能编码记忆，却面临长序列和 context bottleneck。}{官方视频 03:45。}

\begin{knowledgebox}{读图：不要把“RNN 不工作”当绝对结论}
Slide 的红色叉号强调长上下文与内容访问困难，不表示 RNN 在所有任务上失效。门控 RNN 曾经并仍可在小模型、流式处理和强顺序偏置任务中有效。Transformer 的优势来自更短的信息路径和并行训练，但也付出 attention matrix、内存和较弱局部偏置的代价。正确比较应指定序列长度、延迟、数据规模和硬件。
\end{knowledgebox}

\begin{importantbox}{Bottleneck 的本质}
若 decoder 第 $t$ 步只能读取同一个 $c$，它无法直接问“源序列中哪个位置与当前生成最相关”。Attention 的关键不是把向量做得更大，而是把接口改成 query-conditioned retrieval：每个生成步都能产生自己的权重并读取不同位置。
\end{importantbox}

\subsection{本章小结}

Transformer 的前史是一条信息访问路线：RNN 递归携带状态，seq2seq 压缩序列，attention 按需读取，self-attention 则让每个 token 都能成为查询者和信息源。

